% !TEX root = ../../../main.tex
% Sources: LEC02, IMG_9119--IMG_9121.
\section{Fields and order}
\label{sec:analysis-i:fields-and-order}

Before asking which limits exist, we separate the algebra from the order.
A set acquires additional structure when operations are specified.
A \emph{group} is a set with an associative operation, an identity, and an
inverse for each element. It is \emph{abelian} when the operation commutes.

\begin{definition}
  A \emph{field} $K$ has two distinct elements $0$ and $1$, and operations
  of addition and multiplication, such that $(K,+)$ and $(K\setminus\{0\},\cdot)$
  are abelian groups and multiplication distributes over addition.
\end{definition}
These are structures on sets, rather than a literal containment chain of
``sets inside groups inside fields.'' The examples $\Q$ and $\R$ have their
usual arithmetic operations.

The \emph{characteristic} of a field is the least positive integer $n$ such
that $n\cdot1=0$, if there is one, and is zero otherwise. Here $n\cdot1$
means adding $1$ to itself $n$ times.

\begin{proposition}
  Every field of characteristic zero contains a unique smallest subfield,
  canonically isomorphic to $\Q$.
\end{proposition}
\begin{proof}
  Send an integer $m$ to $m\cdot1$, with the usual interpretation for
  zero and negative integers. Characteristic zero makes this map injective.
  For $m\in\Z$ and $n\in\N$, define
  \[
    \iota(m/n)=(m\cdot1)(n\cdot1)^{-1}.
  \]
  The denominator is nonzero. If $m/n=p/q$, then $mq=np$, and the same
  equality holds after mapping the integers into $K$. Multiplying by the
  inverses of the denominators shows that the definition is independent
  of the fraction chosen. The reverse argument shows injectivity.
  The usual common-denominator and product calculations show that $\iota$
  preserves addition and multiplication, so its image is a subfield.
  Any subfield of $K$ contains $1$, its integer multiples, their negatives,
  and the inverses of its nonzero elements. It therefore contains every
  $\iota(m/n)$, proving minimality and uniqueness.
\end{proof}

\subsection{Ordering the arithmetic}
A partial order is a reflexive, antisymmetric, transitive relation $\leq$.
It is a \emph{total order} if any two elements are comparable.
For example, coordinatewise order on $\Q^2$ is partial, but not total;
$(0,1)$ and $(1,0)$ are incomparable. The usual order on $\Q$ is total.

\begin{definition}
  An \emph{ordered field} is a field with a total order satisfying
  \[
    a<b\ \Longrightarrow\ a+c<b+c,
    \qquad
    a<b,\ c>0\ \Longrightarrow\ ac<bc.
  \]
\end{definition}
Every nonzero square is positive, so $1>0$. Consequently every positive
integer multiple of $1$ is positive, and an ordered field has characteristic
zero. It therefore contains the rational subfield just constructed.
The additional property distinguishing $\R$ from $\Q$ is completeness.
